\documentclass{article}
\usepackage{pathfinder-readable}

\title{Truthful Peer Reports and an Uncertain Ranking}

\begin{document}
\maketitle

\begin{abstract}
This note joins Q's method for rewarding truthful reports with P's method for learning a ranking from peer
comparisons. We asked whether truthful peer reports identify the intended ranking. On a four-item path, the
same reports and the same strict incentives occur with opposite endpoint rankings. Under positive noise that
may vary by comparison, the reported edge signs identify exactly the orders joined by directed paths. The
verifier left the thread at DRAFT because these restricted results support a paper section, while the
strategic and decision-level questions remain open.
\end{abstract}

\section{Introduction}

Q, \emph{Mechanism Design for Alignment and Control}, studies how a designer can steer AI agents whose
preferences, abilities, and information are uncertain \cite{Q}. One part uses an agent's prediction of another
agent's report to reward truthful reporting. When different types make different predictions, the score gives
a strict reason to report one's type truthfully. Q also studies action rules, but its general implementation
result assumes a known model of types and beliefs and can use rewards that depend on the true state.

P, \emph{When Peer Agreement Hides a Reversed Ranking}, studies a ranking inferred from several observers'
labels of the same comparisons \cite{P}. Each item has a Bradley--Terry score: the score difference sets the
chance that one item beats another. Under P's clean model, observers make independent errors around each
realised comparison. Their agreement helps calibrate observer reliability, after which their average labels
recover score differences. P also shows that a shared error can leave the whole report distribution unchanged
while reversing a ranking.

We pursued the meeting point of these claims. Could Q's strict incentive to tell the truth make P's intended
ranking trustworthy? The answer depends on what ``truth'' means: an agent can faithfully report the label it
saw even when that label does not identify the comparison the designer cares about. P uses the letters Q and P
for a different internal paper pair; here they always mean the two papers just introduced.

\section{What was done}

The peers first separated the incentive question from the identification question. Q's quadratic score
compares a reported type with that type's known prediction of other reports. P's path example compares two
possible worlds with the same complete report law. The peers then placed Q's score on the observers in that
example. They checked that each observer's private vector of labels predicts the other observer's vector
differently from every alternative vector. The same score table therefore rewards truthful labels in both
worlds.

Next they asked what extra comparisons could expose the hidden ranking. Under P's assumption that the shared
error rate is the same on every comparison, a suitable short cycle identifies that rate. Adding the comparison
between items 1 and 3 to the original path creates such a triangle. The peers then tested the assumption
itself. They found two triangle-and-tail worlds with the same reports and opposite orders for items 1 and 4
when the positive error rate may vary by edge. This led to a more general rule based on the directions of
observed edge signs. Finally, they checked a sample bound for learning those directions from truthful reports.

\section{Results}

\subsection{Truthful reports can conceal an order}

Consider items 1, 2, 3, and 4 compared along a path. Write \(c_v\) for item \(v\)'s intended Bradley--Terry
score and \(d_{uv}=c_u-c_v\) for a score difference. In the shared-error world, the three edge differences are
\[
(d_{12},d_{23},d_{34})=(4,-3/2,-3/2),
\qquad c_1-c_4=1.
\]
Each intended comparison is independently flipped for both observers with a common positive flip mean
\(b=1/2\). In a clean world without that shared flip, use edge differences
\[
f_b(d)=2\operatorname{artanh}\bigl(b\tanh(d/2)\bigr).
\]
The clean world's endpoint difference is \(f_{1/2}(4)-2f_{1/2}(3/2)\simeq-0.26458\). A path permits both sets
of edge differences as score vectors. Its shared corrupted comparison has the same binary distribution on each
edge in the two worlds. Independent draws across edges and the same personal observer errors therefore give
the same \emph{entire} joint distribution of reports, although the intended endpoint orders differ \cite{P}.

The peers' direct check in \url{../ada/truthful_peer_scores_hidden_target.md} supplied the incentive step. Let
\(Y_{ei}\) be observer \(i\)'s label on edge \(e\), let \(a_i\in(0,1)\) be that observer's personal
reliability, and let \(w_e\) be the mean sign of the shared corrupted comparison. For a second observer \(j\),
they obtained
\[
\mathbb E[Y_{ej}\mid Y_{ei}=+1]
-\mathbb E[Y_{ej}\mid Y_{ei}=-1]
=\frac{2a_i a_j(1-w_e^2)}{1-a_i^2w_e^2}>0.
\]
Thus any two distinct private label vectors predict different distributions of the other observer's vector.
Call that prediction \(\beta_i[t]\) for private vector \(t\). Q's quadratic score, scaled by any
\(\Lambda>0\), gives a truthful report an expected advantage \(\Lambda\|\beta_i[t]-\beta_i[s]\|_2^2>0\) over a
false report \(s\), provided the other observer reports truthfully \cite{Q}. The check used a single feasible
action and no other utility. It established a strict truthful equilibrium in both worlds, not a unique
equilibrium or Q's full action-implementation claim.

Any rule that sees only these reports has the same output law in both worlds. Its probabilities of choosing
the wrong endpoint order add to at least one. Hence its error is at least \(1/2\) in one world, however many
independent repetitions it sees.

\subsection{Which orders the reports can certify}

Now allow a separate unknown positive shared-flip mean \(b_e\in(0,1]\) on each comparison edge \(e=(u,v)\).
Assume independent edge draws and positive observer reliabilities. The calibrated mean corrupted sign is
\[
w_e=b_e\tanh\bigl((c_u-c_v)/2\bigr).
\]
For a nonzero \(w_e\), its sign tells which endpoint has the higher intended score. Direct every edge from
higher to lower. The peers proved that a strict order between two items is fixed by the full report law
exactly when a directed path joins them. A path forces each score to fall. If neither item reaches the other,
either order extends the observed directions. In each extension, sufficiently wide score gaps permit positive
values of \(b_e\) that reproduce the same \(w_e\), and hence the same report law. This is the reachability
argument checked by both peers in the ledger and set out in
\url{../ada/truthful_peer_scores_hidden_target.md}. A zero-mean edge instead forces equal endpoint scores
under the positive-flip assumption; feasible ties may be contracted first.

This result is about the order of a particular pair under the stated noise model. Across all possible score
patterns, certifying every pair requires a complete comparison graph. A sparse graph can still certify pairs
linked by directed paths. For a direct comparison, even the numerical reliability need not be known if its
sign is known to be positive.

There is also a conditional sample statement. Suppose every observer used on an edge has reliability at least
\(\alpha>0\), every \(|w_e|\) is at least \(\kappa>0\), and \(N\) independent truthful labels are collected
per edge. For \(|E|\) edges and desired failure probability \(\zeta\), the peers checked that
\[
N\geq\frac{2}{\alpha^2\kappa^2}
\log\frac{2|E|}{\zeta}
\]
makes all edge directions correct with probability at least \(1-\zeta\). Here \(E\) is the comparison-edge
set. Hoeffding's inequality bounds each sign error and a union bound covers all edges. The guarantee then
applies to every order certified by reachability. It assumes truthful reports and a positive signal margin; it
does not show that a score table estimated from strategic reports preserves truthful play.

\section{What did not hold}

The path example does not contradict Q. Q's full type includes beliefs about the state, and its designer is
given a state-and-type model. Its general reward can cancel reported utility and penalise unsupported actions
using the true state. The peers' example deliberately used label vectors as types, a report-only score, and a
single action. It tests what labels alone certify when the intended target is to be learned from them.

P's short-cycle repair also has a precise boundary. A nondegenerate triangle or four-cycle can identify one
common positive flip rate from calibrated population means, so a new edge between items 1 and 3 repairs the
original path \emph{within that constant-rate model} \cite{P}. The peers' triangle-and-tail example showed
that the reports cannot verify that the rate really is constant: positive edge-specific rates can give the
same reports while reversing the order of items 1 and 4. A direct comparison of 1 and 4 would certify that
particular order if the shared flip rates and observer reliabilities remain positive. If their signs can be
negative or unknown, even this sign argument fails.

Peer prediction and the problem of uninformative equilibria were already studied by Miller, Resnick, and
Zeckhauser \cite{Miller2005} and by Shnayder and colleagues \cite{Shnayder2016}. The peers did not establish
novelty for the present composition or graph criterion. They also left equilibrium selection, score tables
learned from strategic reports, stable cycle calibration in finite samples, and a specified downstream
decision payoff unresolved.

\section{Why the thread stopped}

The verifier marked the consolidated note DRAFT. Its stated reason was that the strict-scoring construction,
the identical-report rank reversal, and the conditional reachability result were sound and useful as a paper
section, with their limits stated. They did not yet amount to a finished account of strategic reporting and a
decision benefit. The smallest concrete restart for this pair is to fix the decision between items 1 and 4,
add their direct comparison, and derive its label count and decision-error bound under truthful reporting and
positive edge-specific flips. That would answer this one ordinal audit question; the broader incentive and
noise questions would still require work.

\bibliographystyle{plainurl}
\bibliography{references}

\end{document}
